% ============================================================
% preambulo.tex — Estilo geral da apostila
% ============================================================

% ---- Fontes -------------------------------------------------
\usepackage{fontspec}
\setmainfont{Fira Sans}
\setsansfont{Fira Sans}
\setmonofont{Fira Mono}[Scale=0.88]

% ---- Idioma ---------------------------------------------------
\usepackage{polyglossia}
\setmainlanguage{brazilian}
\AtBeginDocument{
    \renewcommand{\tablename}{Tabela}
    \renewcommand{\figurename}{Figura}
    \renewcommand{\contentsname}{Sumário}
    \renewcommand{\appendixname}{Apêndice}
    \renewcommand{\chaptername}{Capítulo}
    \renewcommand{\partname}{Parte}
}

% ---- Layout de página ----------------------------------------
\usepackage[a4paper,margin=2.5cm,marginparwidth=2.2cm,marginparsep=0.4cm]{geometry}
\usepackage{setspace}
\onehalfspacing

% ---- Cores ------------------------------------------------
\usepackage{xcolor}
\definecolor{primaria}{HTML}{0B3D5C}      % azul petróleo escuro (título)
\definecolor{secundaria}{HTML}{12809E}    % azul-esverdeado (ESP32/PCB)
\definecolor{acento}{HTML}{E8762C}        % laranja circuito (destaques)
\definecolor{fundoclaro}{HTML}{F4F7F8}
\definecolor{fundocodigo}{HTML}{282C34}
\definecolor{textoclaro}{HTML}{E8EEF1}
\definecolor{cinzaesc}{HTML}{3B4750}
\definecolor{notaBg}{HTML}{EAF4F6}
\definecolor{notaBorda}{HTML}{12809E}
\definecolor{atencaoBg}{HTML}{FCEFE6}
\definecolor{atencaoBorda}{HTML}{E8762C}
\definecolor{dicaBg}{HTML}{EEF6EC}
\definecolor{dicaBorda}{HTML}{4C8C4A}

% ---- Matemática / símbolos ---------------------------------
\usepackage{amsmath,amssymb}

% ---- Gráficos e diagramas -----------------------------------
\usepackage{graphicx}
\usepackage{tikz}
\usetikzlibrary{shapes,arrows.meta,positioning,fit,backgrounds,calc}

% ---- Tabelas --------------------------------------------------
\usepackage{array,tabularx,booktabs,longtable,multirow}
\renewcommand{\arraystretch}{1.25}

% ---- Listas -----------------------------------------------
\usepackage{enumitem}
\setlist{itemsep=0.15em, topsep=0.4em}

% ---- Hyperlinks e referências -------------------------------
\usepackage{hyperref}
\hypersetup{
    colorlinks=true,
    linkcolor=primaria,
    citecolor=primaria,
    urlcolor=secundaria,
    pdftitle={Introdução à Linguagem C para Sistemas Embarcados},
    pdfauthor={Apostila EFB2006},
    pdfsubject={Linguagem C, ESP32, PlatformIO},
    bookmarksnumbered=true,
    bookmarksopen=true
}
\usepackage[capitalize,noabbrev]{cleveref}

% ---- Código-fonte (minted + tcolorbox) ------------------------
\usepackage{fontawesome5}
\usepackage{tcolorbox}
\tcbuselibrary{minted,breakable,skins}
\usemintedstyle{one-dark}

\tcbset{
  codigobase/.style={
    breakable, enhanced, listing only, minted style=one-dark,
    minted options={fontsize=\small, breaklines, breakanywhere,
                     tabsize=4, baselinestretch=1.05},
    colback=fundocodigo, colframe=fundocodigo,
    coltitle=textoclaro, fonttitle=\bfseries\ttfamily\small,
    boxrule=0pt, top=5mm, bottom=1mm, left=1mm, right=1mm,
    arc=2mm, outer arc=2mm,
    before={\color{textoclaro}},
    attach boxed title to top left={xshift=3mm, yshift=-1mm},
    boxed title style={colback=acento, arc=1.5mm, boxrule=0pt,
                        top=0.6mm, bottom=0.6mm, left=2mm, right=2mm},
  },
  terminalbase/.style={
    breakable, enhanced, listing only, minted style=one-dark,
    minted options={fontsize=\small, breaklines, breakanywhere,
                     tabsize=4, baselinestretch=1.05},
    colback=fundocodigo, colframe=cinzaesc,
    coltitle=textoclaro, fonttitle=\bfseries\ttfamily\small,
    boxrule=0.4pt, top=1mm, bottom=1mm, left=1mm, right=1mm,
    arc=2mm, outer arc=2mm,
    before={\color{textoclaro}},
  }
}

% Ambientes de código com título automático de arquivo
\newtcblisting{codigoc}[1]{codigobase, minted language=c, title={#1}}
\newtcblisting{codigoini}[1]{codigobase, minted language=ini, title={#1}}
\newtcblisting{codigobash}[1]{terminalbase, minted language=bash, title={#1}}

% ---- Caixas de destaque (nota, atenção, dica) -----------------
\newtcolorbox{nota}{
    breakable, enhanced, colback=notaBg, colframe=notaBorda,
    boxrule=0pt, leftrule=2.6pt, arc=1mm,
    left=3mm, right=3mm, top=1.5mm, bottom=1.5mm,
    fonttitle=\bfseries\color{notaBorda},
    title={\faInfoCircle\ \ Nota}
}
\newtcolorbox{atencao}{
    breakable, enhanced, colback=atencaoBg, colframe=atencaoBorda,
    boxrule=0pt, leftrule=2.6pt, arc=1mm,
    left=3mm, right=3mm, top=1.5mm, bottom=1.5mm,
    fonttitle=\bfseries\color{atencaoBorda},
    title={\faExclamationTriangle\ \ Atenção}
}
\newtcolorbox{dica}{
    breakable, enhanced, colback=dicaBg, colframe=dicaBorda,
    boxrule=0pt, leftrule=2.6pt, arc=1mm,
    left=3mm, right=3mm, top=1.5mm, bottom=1.5mm,
    fonttitle=\bfseries\color{dicaBorda},
    title={\faLightbulb\ \ Dica}
}
\newtcolorbox{esp32box}{
    breakable, enhanced, colback=fundoclaro, colframe=secundaria,
    boxrule=0pt, leftrule=2.6pt, arc=1mm,
    left=3mm, right=3mm, top=1.5mm, bottom=1.5mm,
    fonttitle=\bfseries\color{secundaria},
    title={\faMicrochip\ \ Aplicação no ESP32}
}

% ---- Títulos de capítulo/seção ---------------------------------
\usepackage{titlesec}
\titleformat{\chapter}[display]
  {\normalfont}
  {\Large\color{secundaria}\fontsize{14}{16}\selectfont\textsc{Capítulo \thechapter}}
  {6pt}
  {\Huge\bfseries\color{primaria}}
\titlespacing*{\chapter}{0pt}{-10pt}{28pt}

\titleformat{\section}
  {\normalfont\Large\bfseries\color{primaria}}
  {\thesection}{0.6em}{}
\titleformat{\subsection}
  {\normalfont\large\bfseries\color{secundaria}}
  {\thesubsection}{0.6em}{}
\titleformat{\subsubsection}
  {\normalfont\normalsize\bfseries\color{cinzaesc}}
  {\thesubsubsection}{0.6em}{}

% ---- Cabeçalhos e rodapés -------------------------------------
\usepackage{fancyhdr}
\pagestyle{fancy}
\fancyhf{}
\renewcommand{\headrulewidth}{0.4pt}
\renewcommand{\footrulewidth}{0pt}
\fancyhead[LE,RO]{\small\thepage}
\fancyhead[RE]{\small\color{cinzaesc}\nouppercase{\leftmark}}
\fancyhead[LO]{\small\color{cinzaesc}\nouppercase{\rightmark}}

\fancypagestyle{plain}{
  \fancyhf{}
  \fancyhead[LE,RO]{\small\thepage}
  \renewcommand{\headrulewidth}{0pt}
}

% Páginas em branco (inseridas para iniciar capítulos em página ímpar)
% não devem exibir cabeçalho.
\makeatletter
\renewcommand{\cleardoublepage}{%
  \clearpage
  \if@twoside
    \ifodd\c@page\else
      \hbox{}\thispagestyle{empty}\newpage
      \if@twocolumn\hbox{}\newpage\fi
    \fi
  \fi
}
\makeatother

% ---- Sumário -----------------------------------------------
\usepackage{tocloft}
\renewcommand{\cftchapfont}{\bfseries\color{primaria}}
\renewcommand{\cftchappagefont}{\color{primaria}}
\renewcommand{\cftsecfont}{\color{cinzaesc}}

% ---- Numeração de código e legendas ---------------------------
\usepackage{caption}
\captionsetup{font=small,labelfont={bf,color=primaria}}

% ---- Comando utilitário para inline code ------------------
\newcommand{\code}[1]{\texttt{\color{secundaria}#1}}
\newcommand{\ESP}{ESP32\xspace}
\usepackage{xspace}
